\section{Question and contribution}
Studying how a model acts under incomplete knowledge requires connecting what
it records to what it chooses. Our broader research objective is to understand
how models formulate assumptions, select informative experiments and revise
predictions. Before using external records as evidence about that process,
we ask what those records actually establish: are they available, do their
declared rules reproduce the proposal, and does substituting a rule or pairing
records expose a distinction unavailable from simpler checks? These are
measurement questions. The present experiments do not test whether a model's
chosen experiments resolve more uncertainty or improve held-out predictions.

We study external decision records consisting of a representation and a rule.
In the common-menu experiment, a representation is a numeric table of failure
forecasts and a rule selects a column. These numbers are uncalibrated forecasts;
they are not verified world models or a report of all internal reasoning.
An analyst has privileged access to the task evaluator. A resulting action is
labeled preservation (P), loss (L), or unresolved (U) according to the stated
remaining-task feasibility procedure. None of these labels identifies a model's
internal cause merely because it is attached to an executable record.

The diagnostic proposal examined here originated in our own component pilot:
cross a recorded representation with another recorded rule, then use the
resulting action's task label to characterize their contribution. It is not
presented as an established procedure in widespread use. Three concrete
observations constrain that proposal. First, reproducing a native action did
not uniquely fix the analyst's selector completion: five native-compatible
completions gave different crossed outcomes. Second, optional records were
scarce under the frozen adapter, whereas requiring records exposed same-rule
and row degeneracies. Third, the follow-up's required-record/model-proposal
condition produced eight matching records per profile without externally
selecting the action, yet every such record was action-invariant under the
two supported rules. Sol's matching records also included two conditional L
outcomes. Correspondence, rule sensitivity and task feasibility therefore
need separate readouts in this case.

The practical decision is whether a researcher has enough evidence to
interpret a crossing as a diagnostic contrast, or must first test record
availability, selector specification and within-record sensitivity. For
example, our pilot's apparent off-diagonal P outcome becomes L under another
native-compatible completion (Section~\ref{sec:selector-sensitivity}); calling
that a unique component diagnosis would hide an analyst choice. In the later
common-menu panel, the sole P/L-changing rule contrast is already visible
without a partner record. These examples locate unsupported inferences in
this study; we have not measured whether applying the checks improves other
researchers' diagnosis accuracy or future-behavior prediction.

The contribution is the empirical separation of these measurements and an
executable reader that reconstructs them from identity-bound sources. The
full-scan identity is a supporting structural control, not an empirical
discovery. We do not propose a general attribution algorithm, establish that
crossing is useless across contracts, or demonstrate internal faithfulness.

\section{What the comparison measures}
\subsection{An external computation and two different crossings}
Let $A_i$ be record $i$'s table over a common finite menu and let $r_i$
be its declared deterministic rule. Native action $a_i=r_i(A_i)$ is one
column of that menu. The \emph{observed-rule crossing} has four entries
$c_{ij}=r_j(A_i)$ for $i,j\in\{1,2\}$. The separate
\emph{supported-rule-family table} evaluates each $A_i$ under both allowed
rules: uniform mean and worst case, with lowest-index tie breaking.
The family table may contain a rule that neither paired actor used.
It must not be described as a second live observation or an observed model
counterfactual. Task labels for these actions come from the bound menu census.

\subsection{Same-access baselines and a structural identity}
The comparison holds the archived tables, rule definitions, menu and task-label
map fixed. It removes operations rather than withholding information from a
baseline. We evaluate (i) native actions and rule equality, (ii) native actions
and the single-distinct-row condition, and (iii) independent per-record
evaluation of each supported rule. The last baseline never needs a partner's
table. Only after producing its results do we compare them with paired cells.

\paragraph{Proposition (per-record sufficiency under a fixed contract).}
For fixed tables $A_i$, a common menu, and a fixed family $\mathcal R$ of
deterministic rules with specified tie breaking, the collection
$\{(i,r,r(A_i)):r\in\mathcal R\}$ determines every supported-rule-family
cell and every observed-rule cell whose rule belongs to $\mathcal R$.
\emph{Proof.} Each cell is one of those evaluations; no partner table is an
argument of $r(A_i)$. A fixed action-label map also determines its label.
This statement does not establish the validity of that map or identify
an actor's internal computation.

Two weaker checks need fewer rule evaluations. If $r_1=r_2$, then
$c_{i1}=c_{i2}=a_i$. If a table has a single distinct row, mean and worst case
select the same column. Duplicate-containing tables with other distinct
rows do not satisfy this condition. The full-scan equality is thus a
compositional identity of the contract. The
empirical question is how often the weaker checks suffice and what contrasts
remain; agreement with the full scan is not an independent discovery of the
identity. We measure recoverable cells, not a wall-clock speedup or causal
attribution accuracy.

\paragraph{Task and record at a glance.}
In the continuous task, an actor sees a state/action history and chooses a
bounded two-dimensional control; the remaining objective is to reach and
maintain the target region without violating a path constraint. In the route
task, it sees a discrete history and chooses a legal route action under the
remaining task constraints. In this panel each task supplies a finite action
menu and permits up to three analyses of the observed history. A decision
record is a table: columns are menu actions, rows are uncalibrated candidate
failure forecasts in $[0,1]$, and the declared rule minimizes the mean or
the worst row value. In the older 32-trial panel, the external selector executes
that choice in its mandatory condition. Appendix~\ref{app:common-menu-protocol} specifies the
exact histories, menus and six-field schema.

For a didactic two-action table
$A=\left(\begin{smallmatrix}0&0.4\\0.7&0.4\end{smallmatrix}\right)$,
mean scores $(0.35,0.4)$ select action 0, whereas worst-case scores
$(0.7,0.4)$ select action 1. This invented arithmetic example is not a trial
or a task-label claim; it shows what a per-record scan can reveal.

\subsection{Feasibility labels and access}
P requires a successful admissible continuation witness after the action.
L requires the specified loss certificate; absence of a witness alone is not
L. U preserves unresolved numerical cases. Continuous-control labels are
conditional on the disclosed floating-point and bounding premises in the
supplement. The discrete route family's finite computation has a different
verification contract. Both actors and analyst see the legal histories; only
the analyst receives the hidden mechanisms and privileged evaluator during
the original experiment. Reviewer disclosure now permits inspection of these
mechanisms without retroactively changing what participants received.

Record validity, correspondence and rule sensitivity can be checked without
the physical label map. P/L/U comparisons add its assumptions. The exact
arithmetic helper does not prove the uniform library bounds or attest the
historical runtime; those premises remain open. These are distinct evidence
levels when interpreting outcome differences.

\section{Evidence sets and their denominators}
The first three evidence sets below use crossing-pair denominators.
The separate follow-up in Section~\ref{sec:followup} instead compares A/B/C
conditions within eight matched triplets per profile; these condition pairs
are not crossing pairs. Its 48 assignments include two technical losses.
\begin{table}[htbp]
\centering\small
\begin{tabular}{p{0.28\linewidth}rrrr}
\toprule
Evidence set & Trials & Pairs & Eligible records & Eligible pairs\\
\midrule
Pilot (component study) & 8 & 4 & \multicolumn{2}{c}{2 selected component pairs}\\
Prospective transfer & 48 & 24 & 1 & 0\\
Optional record & 16 & 8 & 1 & 0\\
Mandatory record + execution & 16 & 8 & 16 & 8\\
\bottomrule
\end{tabular}
\caption{The three predecessor evidence sets. Pilot component eligibility is a different
construction from the frozen prospective adapter and common-menu contract.
Empty pair denominators leave the crossing effect unassessed.}
\end{table}

The pilot used four fixed histories with two repeats under the requested
gpt-6-astra/max profile. Its component study
focused on two boundary histories, B and D, after outcomes were available.
The transfer panel used 48 live trials: two families, three instances per
family, two histories per instance, the gpt-6-astra/max and
gpt-6.1-sol/max profiles, and two repeats. Under its frozen adapter, a record had to
support the prescribed operation without post hoc completion. Only one trial
qualified, leaving zero of 24 pairs crossable. This is an applicability result,
not evidence that models lack unrecorded representations or that crossing has
zero diagnostic effect. Native transfer labels were 42 P and 6 L. The adapter's mutually exclusive
outcomes were: 24 route trials with no analysis attempt, 18 whose last successful
analysis lacked a table, three whose rule did not reproduce the action, two
with a menu-identity mismatch, and one supported record. These reasons
distinguish absent evidence from an invalid submitted record.
The other 24, continuous-task trials made 72 analysis attempts: 67 succeeded
and five returned errors. Thus not calling an analysis and calling one that
fails are different observations; neither establishes absence of internal
reasoning or inability to analyze.

The 32-trial panel crossed two task families, two purposefully selected
histories per family, two model profiles (gpt-6-astra/max and gpt-6.1-sol/max),
two repeats and two conditions. Each condition has 16 trials and eight pairs.
A crossing pair comprises the two repeats at the same history and profile
within one condition, not a pair of trials from different conditions.
Both conditions used a common finite menu and structured final output.
Appendix~\ref{app:common-menu-protocol} identifies the selected histories,
execution contract, resource limits and recorded usage.
The optional condition permitted an observational record; the mandatory
condition required a record and let an external deterministic executor select
the action. Mandatory record--action agreement is consequently a construction
property, unlike a matching optional observational record. The combined
intervention does not isolate a pure reporting effect.

All 32 trials were evaluated: no unstarted slots, technical missingness or
protocol-invalid trials remain. This does not mean the analysis interface was
error-free: 48 attempts yielded 41 successes and seven code errors in seven
trials (five optional, two mandatory). Five errors used attribute access on
the dictionary history; two used the unavailable built-in \texttt{next}.
The prompt described \texttt{history.states}/\texttt{history.actions}, although
the worker passed a dictionary, and did not enumerate the built-ins.
Each affected trial consumed all three analysis slots, one on the error and
two successfully. Appendix~\ref{app:interface-audit} gives the trial-level
accounting. No trial is removed or relabeled. The effect on subsequent
actions and record structure is unassessed.
Optional P/L/U counts were 13/3/0; mandatory
counts were 10/5/1. Four purposefully chosen histories do not form a population
sample of tasks or support a model ranking. Table~\ref{tab:label-contracts}
separates the families' label contracts. The study inventory in
Appendix~\ref{app:inventory} identifies excluded historical and development
work; the later acquisition experiment has its own denominator.

\begin{table}[htbp]
\centering\small
\begin{tabular}{llrrrl}\toprule
Condition & Family & P & L & U & Verification contract\\\midrule
Optional & Continuous & 5 & 3 & 0 & Numerical, conditional\\
Mandatory & Continuous & 2 & 5 & 1 & Numerical, conditional\\
Optional & Route & 8 & 0 & 0 & Finite integer recurrence\\
Mandatory & Route & 8 & 0 & 0 & Finite integer recurrence\\\bottomrule
\end{tabular}
\caption{Native outcomes in the common-menu panel, eight trials per row.
Continuous P uses an admissible continuation witness; continuous L depends
on the numerical and historical-runtime premises in Appendix F. U remains
unresolved. Successful present-host replay does not close those premises.}
\label{tab:label-contracts}
\end{table}

\section{Results: applicability is not incremental information}
The mandatory condition made all eight pairs executable. It did not make
their crossing tables informative by construction. Every pair used the same
declared rule, so all 32 observed-rule cells were already determined by native
actions and rule identity. Among the 16 mandatory records, 13 had one row,
one contained duplicate rows alongside distinct rows, and 13 had only one
distinct row. The last count must not be increased by treating every duplicate
as a degenerate table.

\begin{table}[h]
\centering\small
\begin{tabular}{p{0.52\linewidth}rr}
\toprule
Baseline information & Observed cells & Family cells\\
\midrule
Native actions + identical observed rules & 32/32 & ---\\
Native actions + one-distinct-row check & --- & 29/32\\
Independent two-rule scan per record & 32/32 & 32/32\\
Additional cells from full table after scan & 0 & 0\\
\bottomrule
\end{tabular}
\caption{Mandatory condition only. Counts describe recoverable displayed
cells, not independent observations or accuracy on external diagnoses.}
\end{table}

Fifteen of 16 mandatory records selected the same action under both supported
rules. The remaining record changed action and its conditional continuous
P/L category; the independent
scan recovered this change. One pair involved U and remains unresolved rather
than being converted into a resolved loss. The two 32-cell displays overlap;
after within-pair deduplication they reference 30 unique numerical computations,
not 64 new samples. Cross-action labels reuse the existing bound menu census.
No new live model counterfactuals were collected for this retrospective analysis.

\paragraph{Worked same-record contrast.}
For the mandatory \texttt{probe-01} record from gpt-6-astra, repeat 1,
the independent scan yields:
\begin{center}
\begin{tabular}{lrl}\toprule
Rule applied to this record & Menu index & Bound task label\\\midrule
Uniform mean & 7 & L (conditional continuous certificate)\\
Worst case & 12 & P (continuation witness)\\\bottomrule
\end{tabular}
\end{center}
Indices are zero-based. The recorded table, menu identity and label census
are held fixed. No partner record is needed to expose this sensitivity;
it concerns the declared external rule, not a demonstrated internal cause.
The alternative action's label is a deterministic menu evaluation, not an
additional live model decision. This is the sole rule-sensitive mandatory
record, so it must not be presented as a prevalent population effect.
It is also one of the trials affected by the dictionary-access error.
The archived table's rule sensitivity is reproducible; whether correcting
the interface would have changed that table or the decision is unknown.

The optional condition contains one submitted record, which is eligible,
and 15 trials with no submitted record (\texttt{no\_record}). It has no
eligible pairs. Thus 15 absent records must not be described as 15 failed
record implementations. Its incremental effect is unassessed, not zero. Thus this experiment
shows both an applicability bottleneck without mandatory execution and a
redundancy bottleneck after enforcing it. These are distinct conclusions.

\subsection{Pilot contrast and selector sensitivity}
\label{sec:selector-sensitivity}
Each selected pilot history already had one native P and one native L outcome.
The declared B completion produced two P off-diagonals; D produced two L
off-diagonals. These patterns distinguish conditional component dependencies
but introduce no outcome category absent from the native pair. They are not
an independently labeled ground truth for internal model causes.

In B, the off-diagonal operation required an analyst-specified completion.
All five examined completions reproduced the native action, but selecting
the plan with the smallest predicted trajectory peak changed the crossed
outcome from P to L. Appendix~\ref{app:b-completions} defines all five
selectors and reports their outcomes. The conclusion must therefore name
the completion it uses. D's off-diagonal losses show a conditional joint-change pattern, not
that each native decision was a failure. The pilot's heterogeneous optimizers
do not implement the later common-menu comparison; its full-crossing versus
per-record-scan result was not measured. We retain this distinction instead of
extending the later structural identity to an untested setting.


\section{Follow-up: required records with and without external execution}
\label{sec:followup}
The earlier mandatory condition combined a record requirement with an external
executor. The follow-up adds a required-record condition that executes the
model's proposal. Each matched triplet uses one history, repeat and
model/max/Fast profile: A permits an optional record and executes the proposal;
B requires a valid record and executes the proposal; C requires a valid record
and executes its declared selector, retaining the proposal for observation.
Every trial starts a fresh context. Four purposefully selected historical
contexts, two profiles and two repeats give 48 planned assignments, eight
triplets per profile. These repeats are not independent draws of new worlds.
The allocation and six possible within-triplet orders were fixed before
collection; earlier-panel trials are not reused as controls.
Appendix~\ref{app:followup} gives the execution segments, interface and reader
boundaries. No new population sample or live crossed decision was collected.

\begin{table}[htbp]
\centering\small
\begin{tabular}{llrrrrrrr}\toprule
Profile & Arm & Done & Tech. & P/L/U & Valid & Match & One row & Same action\\\midrule
Astra & A & 7 & 1 & 6/1/0 & 0 & -- & 0 & 0\\
Astra & B & 8 & 0 & 8/0/0 & 8 & 8/8 & 8 & 8\\
Astra & C & 7 & 1 & 5/2/0 & 7 & 7/7 & 6 & 6\\\midrule
Sol & A & 8 & 0 & 7/1/0 & 0 & -- & 0 & 0\\
Sol & B & 8 & 0 & 6/2/0 & 8 & 8/8 & 6 & 8\\
Sol & C & 8 & 0 & 7/1/0 & 8 & 8/8 & 7 & 7\\\bottomrule
\end{tabular}
\caption{Follow-up, eight planned assignments per row; profiles are separate
strata, not ranked. Done means a completed trial, not full task success.
Valid counts usable records; Match is record-selected action versus model
proposal among valid records. The empty A denominator is undefined.
Same action counts records for which both supported rules select the same
action. Sol C has one duplicate-containing record; all other rows have zero.
There are no protocol losses, active or unstarted assignments.}
\label{tab:followup-records}
\end{table}

\paragraph{Availability and correspondence.}
In A, no valid record was available among Astra's seven or Sol's eight
completed trials. This is absence of the specified external record, not
absence of reasoning. In B, each profile supplied eight valid records and all
eight matched its own proposal. Unlike C's executed-action agreement, this
correspondence was not imposed by the executor. It does not establish that
the record caused the proposal or faithfully represents internal reasoning.
All eight B records per profile were invariant in action across the two
supported rules. Thus the earlier degeneracy observation also occurs without
external action selection. Single-row and duplicate records were permitted;
their presence is a scientific observation, not a protocol violation.

\paragraph{Conditional outcomes and observed opportunities.}
Table~\ref{tab:followup-pairs} reports the two predeclared contrasts using only
observed nontechnical pairs. Astra has seven fully completed triplets and Sol
eight. L outcomes remain in the comparison; technical failures are reported
separately and are not converted into L. The directions do not support a
common improvement claim across these profiles, nor an equivalence or
no-effect conclusion. In Sol B, eight matching records coexist with two L
labels, illustrating why correspondence and continuation feasibility require
separate measurements. P certifies a conditional feasible continuation;
the experiments do not measure completed-task success.

% Generated from identity-bound existing follow-up sources.
\begin{table}[htbp]
\centering\small
\begin{tabular}{llrrrrrr}\toprule
Profile & Contrast & Pairs & Missing & L$\to$P & P$\to$L & Same & Net P\\\midrule
Astra & B--A & 7/8 & 1 & 1 & 0 & 6 & +1\\
Astra & C--B & 7/8 & 1 & 0 & 2 & 5 & -2\\
Sol & B--A & 8/8 & 0 & 0 & 1 & 7 & -1\\
Sol & C--B & 8/8 & 0 & 2 & 1 & 5 & +1\\
\bottomrule
\end{tabular}
\caption{Descriptive paired transitions run from A to B for B--A, and from B to C for C--B. Pairs are observed/planned; Missing counts technical missing pairs; Same includes P-to-P and L-to-L. Dividing Net P by observed pairs recovers the reported paired difference. No pooled estimate or population inference. C--B is the total condition contrast, not an isolated executor effect on the same fixed record.}
\label{tab:followup-pairs}
\end{table}


\paragraph{Where the outcome contrasts occur.}
The family breakdown in Table~\ref{tab:followup-families} locates every
observed P/L movement in the two continuous-control histories. All completed
route decisions were P, although each route menu contained an L option;
this pattern was not forced by having only P choices. The continuous menus
have 13 actions and the route menus three, so these purposeful contexts are
not a calibrated difficulty ladder. Appendix~\ref{app:followup-history}
reports every profile/history/arm cell. In particular, Sol's C--B net gain
of one P comprises two L-to-P and one P-to-L transitions. Net change alone
does not describe those opposing movements, and the underlying continuous
labels retain the numerical/runtime premises stated above.

% Generated from identity-bound existing follow-up sources.
\begin{table}[htbp]
\centering\small
\begin{tabular}{llrrrr}\toprule
Profile & Family & Assigned & Done & Tech. & P/L/U\\\midrule
Astra & Continuous & 12 & 12 & 0 & 9/3/0\\
Astra & Route & 12 & 10 & 2 & 10/0/0\\
Sol & Continuous & 12 & 12 & 0 & 8/4/0\\
Sol & Route & 12 & 12 & 0 & 12/0/0\\
\bottomrule
\end{tabular}
\caption{Follow-up outcomes by family, separately within each profile. Each family has two purposeful histories and two repeats per arm. Done denotes a completed decision, not full task success. P/L labels in the continuous family retain the numerical/runtime premises.}
\label{tab:followup-families}
\end{table}


\paragraph{A changed action need not change the task label.}
The follow-up contains two rule-sensitive records, both in C
(Table~\ref{tab:followup-rule-examples}). Mean and worst case select different
menu indices, while the bound label remains P in one and L in the other.
These are existing-record calculations, not new live counterfactual trials.
Together with the older panel's single label-changing example, they make
action sensitivity and task-label sensitivity separately inspectable.

% Generated from identity-bound existing follow-up sources.
\begin{table}[htbp]
\centering\small
\begin{tabular}{llrlll}\toprule
Profile & History & Repeat & Arm & Mean & Worst case\\\midrule
Astra & probe-01 & 2 & C & 1 (P) & 0 (P)\\
Sol & probe-02 & 1 & C & 11 (L) & 0 (L)\\
\bottomrule
\end{tabular}
\caption{The two rule-sensitive follow-up records. Entries give zero-based menu index and bound label. The action changes but the label does not; these are per-record calculations on existing menus, not live crossed decisions. Continuous labels remain conditional.}
\label{tab:followup-rule-examples}
\end{table}


Every completed C proposal already matched the record-selected action
(Astra 7/7, Sol 8/8). No observed C trial therefore tests the benefit of
overriding a conflicting proposal. The C--B contrast is not an isolated
mechanical correction effect on the same fixed record. Labels and
within-record rule diagnostics are tied to the existing menu contract;
they are not additional live crossing observations or independent evidence
about new task physics.

\section{Position relative to existing work}
CAR estimates effects of interventions on agent steps through same-policy
downstream replay and also studies interactions \citep{shah2026car}.
Our object is a declared external record; a crossed action with a privileged
feasibility label is not a replay estimate of the stochastic policy's effect.
AgentRx uses constraint evidence and a judge to diagnose execution traces
\citep{barke2026agentrx}. AGENTSCOPE combines execution dependencies with
neural invariants for localization \citep{bi2026agentscope}. We do not evaluate
their failure-location accuracy or claim that structured evidence is new.

BLADE is a particularly relevant comparator: its dependency-guided handoff
audits separate information use from the final attribution, and its ablations
compare richer evidence with simpler configurations and selective with
exhaustive inspection \citep{guo2026blade}. Consequently, neither the general
warning that structure alone is insufficient nor the value of a simpler
baseline is our novelty claim. Our specific contribution is the decomposition
of this crossing contract's outputs into applicability, same-rule identity,
row degeneracy and remaining per-record rule sensitivity.

Faithfulness requires evaluating the relation between an explanation and the
system it purports to explain \citep{jacovi2020faithful}. Interchange
interventions on actual neural representations test a different relation
\citep{geiger2021causal}. Our record substitutions neither intervene on those
representations nor establish that the external record generated the model's
internal decision. The literature comparison is scoped to inspected primary
sources; it is not an exhaustive novelty proof or a head-to-head experiment.

\paragraph{Behavior-linked explanation evaluation.}
\citet{madsen2024self} test self-explanation consistency against model behavior
using counterfactual, attribution and redaction tasks. \citet{mayne2026positive}
measure whether adding an explanation helps a predictor anticipate a reference
model's answers on related inputs, relative to the same information without
the explanation. Our record check instead recomputes a declared action from
an archived table. It observes correspondence on that decision; it does not
measure either behavior under the explanation's proposed input changes or
incremental predictive information on new decisions. Those targets provide
concrete extensions, not results established by our matching-record counts.

\paragraph{Connection to scientific discovery.}
BoxingGym evaluates experiment selection through expected information gain
and model discovery through predictions, including whether a scientist's
explanation helps another agent predict without further experiments
\citep{gandhi2025boxinggym}. Our study examines how to interpret a decision
record within such a broader research agenda; it does not benchmark those
discovery abilities or outperform BoxingGym. The distinction is between a
record-level measurement and independently assessed experimental or predictive
utility. A correct external calculation alone cannot supply the latter.

\section{Interpretation and remaining questions}
The case supports a sequence of measurement decisions. An absent record
leaves the proposed operation unsupported. A native-compatible record may
still leave selector completion ambiguous. A matching proposal may be
action-invariant under the supported rules and may receive an L label.
Even a rule-sensitive action may preserve its task label. Each observation
answers a different question; none by itself establishes an internal cause
or a gain in diagnostic accuracy.

For researchers using this contract, these distinctions determine which
additional analysis is justified. The transfer panel's empty pair denominator
calls for an applicability assessment before interpreting crossing effects.
The pilot requires reporting selector-completion sensitivity. In the
common-menu panel, a per-record scan already exposes the remaining rule
sensitivity; the paired display adds no independently evaluated diagnostic
target. The follow-up shows that matching records and action invariance can
coexist without enforced action selection. It does not establish that these
checks improve real-world diagnosis or that a different contract is redundant.

The broader research program remains open. Three next questions connect
records to behavior: whether a stated assumption predicts the next experiment
choice; whether that choice reduces more uncertainty than a valid fixed plan
with the same resources; and whether a revised explanation predicts cases
withheld during its construction. Each needs an independently scored outcome
and a design that gives the behavior a real opportunity to occur. The present
purposeful histories, repeated contexts and fixed two-rule family do not answer
those questions. More trials under the identical crossing contract would not
turn its full-scan identity into evidence of added diagnostic value.

Several boundaries also remain empirical. B correspondence does not identify
whether the record caused the proposal. All completed C proposals already
matched their selectors, leaving mechanical correction of conflicting
proposals unmeasured. The seven analysis errors in the older 32-trial panel
have an unisolated effect; their occurrence includes its sole label-changing
rule-sensitive record. A later corrected interface does not establish what
that earlier record would have been. Continuous task labels retain their
numerical/runtime premises, whereas record-level calculations do not require
those physical claims. P is a continuation witness, not measured full-task
completion. No population effect, profile ranking, calibrated difficulty scale
or internal reasoning faithfulness follows from this study.

An optional companion note addresses one bounded part of the broader program:
selecting an additional observation for target certification in a finite class
of deterministic systems. Its 5,000 paired draws yield a 1.78 percentage-point
gain with interval [1.23, 2.29] points. It uses deterministic policies and
additional local computation, not live LLM discovery. The note's methods and
compute costs are separate supporting material; its result is not needed to
establish the record findings above and is not pooled with them.

\section{Reader access and reproducibility}
\label{sec:reader-access}
The accompanying review bundle contains the LaTeX sources, disclosed
mechanisms, pilot/transfer numerical replay and the 32-trial record comparison.
It also regenerates all 32 menu labels from identity-bound histories and
mechanisms. Exact condition instructions, schema, all 80 archived request
payloads and a renderer recovering their original prompt hashes are included;
archived participant analysis text is never executed.

The companion's offline path reconstructs the original catalog and exclusions,
seeded rejection draw, 5,000 eight-step histories, both policies' queries and
certificates, target audits, counts, work and primary interval. Original
committed record bytes and frozen sources are checked before comparison.
This deterministic replay of existing inputs is distinct from collecting new
observations or reattesting historical processes and elapsed times.

The follow-up supplies all 46 completed minimal final JSON responses,
including the 31 numeric decision tables and their declared rules, the
original allocation, and the common and A/B/C instructions with their wire
schema. The standard-library reader derives record validity, declared-rule
selection, both supported rule actions and row diagnostics directly from
the preserved numeric lexemes. It compares the derived decisions to the
archived results, verifies menu and profile bindings, and checks all 46
completed labels by lookup in the existing bound menu census. This is a
source-table and label-link check, not fresh physical grading or independent
historical attestation. Participant analysis programs, intermediate analysis
outputs and raw trajectories are not included. The older scientific inputs
remain unchanged; this revision did not rerun their physics.

From the bundle root, run:
\begin{verbatim}
python3.12 -I -B -S verify_all.py --followup-only
\end{verbatim}
This
checks the exact package, source-derived follow-up decisions, family/history
counts, directional paired changes and the generated manuscript tables. Omitting the flag also runs
the unchanged older numerical and acquisition checks; the README describes
their previously verified scope and current validation boundary. No model account or service is needed. Same-host or clean-directory
replay is not independent laboratory replication. Historical evidence remains
preserved; this preparation is not a claim of public release, submission or
acceptance.


\paragraph{Responsible interpretation.}
The experiments use simulated environments and model outputs, not human-subject
data. Agreement with an executable record is not evidence of honest intentions,
internal reasoning faithfulness, or safe deployment. These distinctions matter
if the checks are reused to assess consequential agent decisions. The disclosed
mechanisms support inspection of this fixed study, not a claim that a future
benchmark remains secret or resists contamination.
